All code used in this work is available at the project repository.
The data-generation pipeline consists of the following components:
\texttt{src/data/generate\_waveforms.py} for IMRPhenomD waveform
generation; \texttt{src/data/inject\_deviations.py} for amplitude,
phase, and frequency modulation; \texttt{src/data/build\_dataset.py}
for the full data pipeline with whitening, bandpass, and
normalization; \texttt{src/models/cnn1d.py} for the hybrid CNN +
statistics model; \texttt{src/training/train.py} for the training
loop with early stopping and gradient clipping; and
\texttt{scripts/build\_gw150914\_dataset.py} for the GW150914
template study.

Software versions used are Python 3.11, PyTorch 2.x, PyCBC 2.x,
GWpy 3.x, NumPy 1.26, SciPy 1.11+ and scikit-learn 1.3+.

Training used the Adam optimizer with learning rate $3 \times 10^{-3}$,
weight decay $10^{-4}$, batch size $32$ or $64$, cosine annealing of
the learning rate, gradient clipping at norm $1$, early stopping
with patience $20$--$50$ epochs, maximum $200$--$300$ epochs, and
binary cross-entropy loss.

Data were sampled at $\SI{4096}{Hz}$ with a $\SI{4}{s}$ window
($L = 16384$ samples), a low-frequency cutoff of $\SI{20}{Hz}$, and
a bandpass range of $20$--$\SI{500}{Hz}$. Reference masses were
$30 + 25 M_\odot$ unless otherwise stated.

The GW150914 template study used the file
\texttt{H-H1\_GWOSC\_4KHZ\_R1-1126259447-32.hdf5} from the GWOSC
event API for GWTC-1-confident/GW150914/v3. The file contains
$131{,}072$ samples at $\SI{4096}{Hz}$, starting at GPS
$1126259447$. The event at GPS $1126259462.4$ is located at sample
$63078$ of the file.
